% ---- pg-444 （书页 432） ----

挑出的 $n=784$ 个行政人员, 测量了 $p=5$ 个职业特性和 $q=7$ 个职业满意度变量. 职业满意度的测量与职业特性是否相联系? 这个答案可能对职业设计有启示.

原始职业特性变量 $\mathbf{X}^{(1)}$ 和职业满意度变量 $\mathbf{X}^{(2)}$ 分别定义为
\[
\mathbf{X}^{(1)}=\begin{bmatrix}X_1^{(1)}\\X_2^{(1)}\\X_3^{(1)}\\X_4^{(1)}\\X_5^{(1)}\end{bmatrix}=\begin{bmatrix}\text{用户反馈}\\\text{任务重要性}\\\text{任务多样性}\\\text{任务特性}\\\text{自主权}\end{bmatrix}
\]
\[
\mathbf{X}^{(2)}=\begin{bmatrix}X_1^{(2)}\\X_2^{(2)}\\X_3^{(2)}\\X_4^{(2)}\\X_5^{(2)}\\X_6^{(2)}\\X_7^{(2)}\end{bmatrix}=\begin{bmatrix}\text{主管满意度}\\\text{事业前景满意度}\\\text{财政满意度}\\\text{工作强度满意度}\\\text{公司认同感}\\\text{工种满意度}\\\text{总体满意度}\end{bmatrix}
\]
变量 $\mathbf{X}^{(1)}$ 和 $\mathbf{X}^{(2)}$ 应答值在一个尺度下记录且被标准化了. 对 784 个被测者的样本相关系数矩阵为
\[
\mathbf{R}=\left[\begin{array}{c:c}
\mathbf{R}_{11} & \mathbf{R}_{12}\\
\hdashline
\mathbf{R}_{21} & \mathbf{R}_{22}
\end{array}\right]
\]
\[
=\left[\begin{array}{ccccc:ccccccc}
1.0 & & & & & 0.33 & 0.32 & 0.20 & 0.19 & 0.30 & 0.37 & 0.21\\
0.49 & 1.0 & & & & 0.30 & 0.21 & 0.16 & 0.08 & 0.27 & 0.35 & 0.20\\
0.53 & 0.57 & 1.0 & & & 0.31 & 0.23 & 0.14 & 0.07 & 0.24 & 0.37 & 0.18\\
0.49 & 0.46 & 0.48 & 1.0 & & 0.24 & 0.22 & 0.12 & 0.19 & 0.21 & 0.29 & 0.16\\
0.51 & 0.53 & 0.57 & 0.57 & 1.0 & 0.38 & 0.32 & 0.17 & 0.23 & 0.32 & 0.36 & 0.27\\
\hdashline
0.33 & 0.30 & 0.31 & 0.24 & 0.38 & 1.0 & & & & & &\\
0.32 & 0.21 & 0.23 & 0.22 & 0.32 & 0.43 & 1.0 & & & & &\\
0.20 & 0.16 & 0.14 & 0.12 & 0.17 & 0.27 & 0.33 & 1.0 & & & &\\
0.19 & 0.08 & 0.07 & 0.19 & 0.23 & 0.24 & 0.26 & 0.25 & 1.0 & & &\\
0.30 & 0.27 & 0.24 & 0.21 & 0.32 & 0.34 & 0.54 & 0.46 & 0.28 & 1.0 & &\\
0.37 & 0.35 & 0.37 & 0.29 & 0.36 & 0.37 & 0.32 & 0.29 & 0.30 & 0.35 & 1.0 &\\
0.21 & 0.20 & 0.18 & 0.16 & 0.27 & 0.40 & 0.58 & 0.45 & 0.27 & 0.59 & 0.31 & 1.0
\end{array}\right]
\]
$\min(p,q)=\min(5,7)=5$ 个样本典型相关系数和样本典型变量系数向量(取自邓纳姆[4]), 在下表中列出

\begin{mvtable}
\mvtabcap{典型变量系数和典型相关系数}
\mvfit{%
\begin{tabular}{lcccccclccccccc}
\toprule
& \multicolumn{5}{c}{标准化变量} & & \multicolumn{7}{c}{标准化变量}\\
\cmidrule(lr){2-6}\cmidrule(lr){9-15}
& $z_1^{(1)}$ & $z_2^{(1)}$ & $z_3^{(1)}$ & $z_4^{(1)}$ & $z_5^{(1)}$ & $\hat{\rho}_1^*$ & & $z_1^{(2)}$ & $z_2^{(2)}$ & $z_3^{(2)}$ & $z_4^{(2)}$ & $z_5^{(2)}$ & $z_6^{(2)}$ & $z_7^{(2)}$\\
\midrule
$\hat{\mathbf{a}}_1'$: & 0.42 & 0.21 & 0.17 & $-0.02$ & 0.44 & 0.55 & $\hat{\mathbf{b}}_1'$: & 0.42 & 0.22 & $-0.03$ & 0.01 & 0.29 & 0.52 & $-0.12$\\
$\hat{\mathbf{a}}_2'$: & $-0.30$ & 0.65 & 0.85 & $-0.29$ & $-0.81$ & 0.23 & $\hat{\mathbf{b}}_2'$: & 0.03 & $-0.42$ & 0.08 & $-0.91$ & 0.14 & 0.59 & $-0.02$\\
$\hat{\mathbf{a}}_3'$: & $-0.86$ & 0.47 & $-0.19$ & $-0.49$ & 0.95 & 0.12 & $\hat{\mathbf{b}}_3'$: & 0.58 & $-0.76$ & $-0.41$ & $-0.07$ & 0.19 & $-0.43$ & 0.92\\
$\hat{\mathbf{a}}_4'$: & 0.76 & $-0.06$ & $-0.12$ & $-1.14$ & $-0.25$ & 0.08 & $\hat{\mathbf{b}}_4'$: & 0.23 & 0.49 & 0.52 & $-0.47$ & 0.34 & $-0.69$ & $-0.37$\\
$\hat{\mathbf{a}}_5'$: & 0.27 & 1.01 & $-1.04$ & 0.16 & 0.32 & 0.05 & $\hat{\mathbf{b}}_5'$: & $-0.52$ & $-0.63$ & 0.41 & 0.21 & 0.76 & 0.02 & 0.10\\
\bottomrule
\end{tabular}
}
\end{mvtable}
